% Comando simples para exibir comandos Latex no texto
\newcommand{\comando}[1]{\textbf{$\backslash$#1}}

Musical expression is one of humanity’s earliest experiences and one of the most valued forms of art. This art form has been cultivated since ancient times. The Chinese, three thousand years before Christ, had already developed complex systems of musical theory, and Greek culture not only cultivated its study but also associated it with religion, attributing divine characteristics to protect and honor music itself (\textcite{med1996teoria}).

Although music is a subjective form of human knowledge, the recorded music industry today moves approximately US\$29 billion per year (\textcite{ifpi2025}), and thus its value is no longer only subjective.

It is precisely in this context that synthesized sound appears in history. It emerged after the first electronic instruments and, due to the periodic nature of sound, the earliest synthesizers were composed of simple analog oscillators that could be manually controlled through operations such as summing and filtering (\textcite{holmes2016electronic}).

In the 1980s, with the release of the Yamaha DX7 synthesizer, the history of music synthesis took a major turn. The Yamaha DX7 marked the first major success of a revolutionary technique developed by John M. Chowning at Stanford University and later popularized by Yamaha (\textcite{claesson2021resynthesis}). This event represented a new milestone in sound synthesis, not only because of its popularity, but also because it made the reproduction of both realistic instrument simulations and entirely new kinds of sounds accessible to musicians and composers of all levels. In fact, FM synthesis was so successful that the technique remains in use to this day (\textcite{holmes2016electronic}).

This long-standing relevance is especially important because, despite the emergence of newer approaches such as wavetable synthesis, sample-based methods, and GAN-based sound generation, FM synthesis continues to occupy an important place in the contemporary musical ecosystem. Each of these alternatives has strengths and limitations: sample-based methods can be difficult to design and are limited in their ability to generate genuinely new timbres, while GAN-based approaches can produce highly realistic results but often require substantial computational resources. By contrast, FM synthesis remains attractive because it is relatively accessible, flexible, well suited for experimenting with new and even unpredictable sounds, and adaptable to real-time composition.

Thus, the primary goal of the present work is to guide a synthesizer based on frequency modulation (FM) to reproduce instrument sound samples, while also developing a process that fully explores the expressive capabilities of the FM synthesizer itself, achieving an approximation for the resynthesis of arbitrary input sounds.

This approach aligns with one of the key historical advantages of FM synthesizers: their ability to create sounds that cannot be produced by traditional acoustic instruments.

Within this context, researchers continue to focus on improving the control of traditional FM synthesis parameters, which, when done manually, can be a tedious empirical process. For instance, \textcite{claesson2021resynthesis} implemented a neural network to predict FM synthesizer parameters using frequency-domain sound representations as input. Similarly, \textcite{steinmetz2022ddx7} proposed a TCN-based approach that couples the FM synthesizer directly within the network training process.

The present work, however, employs a Convolutional Neural Network (CNN) to achieve its goals, because CNNs enable the direct use of raw audio samples as model input while also serving as effective feature extractors for sound signals (\textcite{kiranyaz2021survey}). CNNs are also well known for their ability to economize computational resources when handling high-dimensional data, which makes them a strong candidate for this type of research task.

Furthermore, beyond a simple resynthesis exercise, this work aims to develop a process that allows the model to achieve a broad understanding of the entire parameter space of the applied FM synthesizer. The idea behind the proposed methodology is that, if the model can learn the FM synthesizer’s parameter space in a comprehensive way, it can then be used to simulate any sound within the synthesizer’s capabilities, making the experience of a future user, such as a musician or composer, much more efficient in achieving the desired sound and integrating it into real composition workflows.

In addition, this work aims not only to apply a machine learning technique, but also to propose an approach that can be adapted to different FM synthesizers, providing practical utility for users interested in sound design and music production.

The remainder of this document is organized as follows. Section~\ref{sec:introduction} presents the fundamental concepts related to music. Section~\ref{sec:related} discusses the related work. Section~\ref{sec:our_proposal} describes the details of the proposed approach and the experimental setup. Section~\ref{sec:experiments} presents and analyzes the results. Finally, Section~\ref{sec:conclusion} provides the conclusions of this work and outlines directions for future research.
